% SCRUM-99 / A2-T1.1 — fragment; requires vntex, longtable, array, hyperref.
% Canonical content: docs/requirements/context-and-scope.md
\section{Bối cảnh, stakeholder và phạm vi hệ thống BrewLite}
\label{sec:context-and-scope}
\subsection{Cơ sở khảo sát và phạm vi đặc tả}

Phạm vi BrewLite được xác định từ Requirements Specification v3.1, stakeholder và các quy tắc BR-01--BR-06. Các yêu cầu dưới đây mô tả bài toán và giới hạn sản phẩm; chức năng, API và quyền người dùng được đối chiếu với implementation ở các phần sau.
\par

Các quy tắc BR-01--BR-06 là cơ sở mô hình hóa vòng đời đơn và thanh toán. Phạm vi đặc tả được phân biệt với chức năng hiện thực và kết quả thử nghiệm.
\par

Nguồn khảo sát gồm Requirements Specification, User Stories, API Contract và tài liệu kiến trúc của dự án; danh mục và phiên bản nguồn được dẫn ở phụ lục\ref{sec:report-provenance}.
\par

\subsection{Bối cảnh và vấn đề}

BrewLite là ứng dụng web đặt cà phê và đồ uống theo mô hình một cửa hàng (Single-Vendor), phục vụ Guest và Customer. Luồng giá trị chính là xem menu, chọn size/topping, quản lý giỏ hàng, đăng nhập, tạo đơn, thanh toán và xem xác nhận/trạng thái/lịch sử đơn.
\par

Bài toán là duy trì thông tin món, cấu hình, giá và kết quả đơn nhất quán xuyên suốt luồng đặt hàng. Giỏ hàng có thể cũ so với dữ liệu cửa hàng; nhiều khách có thể cùng đặt món cuối; request hoặc thông báo thanh toán có thể được gửi lặp hay đến trễ. Hệ thống cần kiểm tra dữ liệu phía server để tránh tính sai tiền, bán vượt tồn kho và lặp tác động thanh toán hoặc cộng điểm.
\par

Phạm vi khảo sát dựa trên đặc tả nghiệp vụ và mô hình một cửa hàng; chưa có khảo sát định lượng về thời gian phục vụ hoặc tỷ lệ sai sót tại một cửa hàng cụ thể.
\par

\subsection{Stakeholder, vai trò và nhu cầu}
\begingroup\small
\renewcommand{\arraystretch}{1.2}
\begin{longtable}{>{\raggedright\arraybackslash}p{0.20\linewidth}>{\raggedright\arraybackslash}p{0.43\linewidth}>{\raggedright\arraybackslash}p{0.25\linewidth}}
\caption{Stakeholder và nhu cầu}\label{tab:auto-context-and-scope-1}\\
\hline
\textbf{Stakeholder} & \textbf{Vai trò và nhu cầu} & \textbf{Giới hạn} \\ \hline
\endfirsthead
\textbf{Stakeholder} & \textbf{Vai trò và nhu cầu} & \textbf{Giới hạn} \\ \hline
\endhead
Guest & Người chưa đăng nhập; cần xem menu, giá và tùy chọn size/topping để lựa chọn món. & Phải đăng nhập trước khi tạo Order. \\ \hline
Customer & Người đặt món; cần giỏ hàng, tạo đơn, thanh toán, xác nhận, lịch sử và điểm tích lũy; cần dữ liệu đúng và bảo vệ quyền sở hữu đơn. & Chỉ truy cập tài nguyên thuộc tài khoản; không quyết định giá, tồn kho hay kết quả thanh toán. \\ \hline
Quản lý cửa hàng / PO mô phỏng & Stakeholder xác định giá trị, ưu tiên backlog và chấp nhận increment; cần phạm vi rõ và quy tắc nhất quán. & Không suy ra có tài khoản hay UI quản trị; vai trò không hàm ý có giao diện quản trị. \\ \hline
Nhân viên quầy / Barista & Stakeholder liên quan đến chuẩn bị và hoàn tất đồ uống; cần hiểu các trạng thái PREPARING, READY, COMPLETED. & UI Barista nằm ngoài MVP; stakeholder nghiệp vụ không đồng nghĩa actor có giao diện đã triển khai. \\ \hline
VNPay Sandbox & Hệ thống ngoài tham gia redirect và trao đổi kết quả IPN/QueryDR; cần dữ liệu giao dịch và thông báo được xác minh theo hợp đồng tích hợp. & Chỉ Sandbox; không giao dịch tiền thật. Mock và ZERO\_AMOUNT là cơ chế nội bộ, không phải actor ngoài. \\ \hline
Nhóm phát triển & Phân tích, thiết kế, triển khai, review và kiểm thử; cần nguồn yêu cầu thống nhất, khả năng truy vết và tái tạo môi trường. & Cần phối hợp đặc tả, thiết kế, source và kiểm thử. \\ \hline
Giảng viên & Đánh giá đồ án; cần tạo phẩm, demo và bằng chứng có thể kiểm tra theo tiêu chí Assignment. & Đối chiếu kết quả với yêu cầu học phần. \\ \hline
\end{longtable}
\endgroup

Guest và Customer là actor sử dụng web chính; VNPay Sandbox là actor hệ thống ngoài. Các stakeholder còn lại có lợi ích và trách nhiệm liên quan nhưng không mặc nhiên có quyền thao tác trong ứng dụng. [S1 §5--6; S2 §0]
\par

\subsection{Mục tiêu sản phẩm}
\begingroup\small
\renewcommand{\arraystretch}{1.2}
\begin{longtable}{>{\raggedright\arraybackslash}p{0.24\linewidth}>{\raggedright\arraybackslash}p{0.68\linewidth}}
\caption{Mục tiêu và kết quả mong đợi}\label{tab:auto-context-and-scope-2}\\
\hline
\textbf{Mục tiêu} & \textbf{Kết quả mong muốn trong baseline} \\ \hline
\endfirsthead
\textbf{Mục tiêu} & \textbf{Kết quả mong muốn trong baseline} \\ \hline
\endhead
G1 -- Đặt món rõ ràng & Khách xem món/giá, chọn size/topping và quản lý giỏ trước khi tạo đơn. \\ \hline
G2 -- Đơn hàng nhất quán & Server kiểm tra giỏ, tính lại giá, lưu thông tin đơn và kiểm soát tồn kho để tránh oversell. \\ \hline
G3 -- Thanh toán có kiểm soát & Kết quả thanh toán được xác minh; request lặp không lặp tác động; retry và trường hợp chưa rõ kết quả tuân thủ baseline. \\ \hline
G4 -- Theo dõi và bảo vệ dữ liệu & Customer xem kết quả/lịch sử đơn của mình; JWT và ownership bảo vệ nghiệp vụ cần đăng nhập; coupon/điểm tuân thủ quy tắc đã có. \\ \hline
G5 -- Có thể đánh giá và bàn giao & Tài liệu mô tả đúng hệ thống, môi trường có thể tái tạo và bằng chứng được ghi theo trạng thái thực tế. \\ \hline
\end{longtable}
\endgroup

Các mục tiêu được cụ thể hóa thành yêu cầu chức năng, phi chức năng và tiêu chí kiểm thử trong SRS.
\par

\subsection{Ranh giới hệ thống}

Bên trong BrewLite: Web Client Next.js, Backend API NestJS, dữ liệu PostgreSQL qua Prisma và các cơ chế nội bộ phục vụ Order, Inventory, Payment, Promotion/Loyalty. Docker Compose hỗ trợ tái tạo môi trường local.
\par

Web Client hiển thị và gửi yêu cầu; Backend là nơi quyết định giá, kiểm tra tồn kho, quyền truy cập và thay đổi trạng thái nghiệp vụ. Không tin giá, userId hay kết quả thanh toán do client tự khai.
\par

Bên ngoài BrewLite: Guest/Customer và VNPay Sandbox. VNPay Sandbox là provider mặc định. Mock CARD/WALLET chỉ phục vụ automated test/demo khi backend cho phép trong môi trường không phải production và luôn bị tắt ở production. ZERO\_AMOUNT xử lý nội bộ đơn 0 VND, độc lập với Mock. Không tự fallback sang Mock khi VNPay PENDING hoặc chưa rõ kết quả. [S1 §3, §9.9, §16]
\par

\subsection{Phạm vi chức năng BrewLite}

Phạm vi gồm hành trình đặt hàng và các cơ chế bảo đảm thanh toán nhất quán. State machine, idempotency, concurrency/inventory và coupon/loyalty là các yêu cầu nghiệp vụ trong phạm vi, không phải tính năng tùy chọn do cách chia bảng. [S1 §4.1, §6; S3 §4.1]
\par

\subsubsection{Đặt hàng, tài khoản và quản lý giao dịch}

\begingroup\small
\renewcommand{\arraystretch}{1.2}
\begin{longtable}{>{\raggedright\arraybackslash}p{0.20\linewidth}>{\raggedright\arraybackslash}p{0.43\linewidth}>{\raggedright\arraybackslash}p{0.25\linewidth}}
\caption{Phạm vi đặt hàng và tài khoản}\label{tab:auto-context-and-scope-3}\\
\hline
\textbf{Nhóm} & \textbf{Nội dung trong phạm vi MVP} & \textbf{Nguồn} \\ \hline
\endfirsthead
\textbf{Nhóm} & \textbf{Nội dung trong phạm vi MVP} & \textbf{Nguồn} \\ \hline
\endhead
Identity & Đăng ký/đăng nhập JWT, băm mật khẩu Argon2id và bảo vệ API/quyền sở hữu dữ liệu. & S1 §4.1, §7.1; S3 §4.1 Task 7 \\ \hline
Catalog và Cart & Menu, chi tiết, size/topping; thêm/sửa/xóa giỏ và tổng tạm; server tính lại khi tạo đơn. & S1 §4.1, §6 F02--F03; S3 §4.1 Task 2--5 \\ \hline
Order và Inventory & Tạo đơn, kiểm tra hợp lệ, snapshot, xác nhận/trạng thái/lịch sử; state machine, chống oversell, concurrency/inventory và hết hạn/hoàn tài nguyên theo lifecycle baseline. & S1 §6 F04--F05; S3 §4.1 Task 6, 9--10 \\ \hline
Mock Payment và idempotency & Thanh toán giả lập CARD/WALLET cho test/demo trong môi trường được backend cho phép; xử lý kết quả và idempotency. Mock luôn bị tắt ở production; policy tại mục 5 áp dụng đầy đủ. & S1 §4.1, §6 F06; S3 §4.1 Task 8, 10 \\ \hline
Promotion và Loyalty & Kiểm tra coupon, giới hạn lượt dùng, tính giảm giá và cộng điểm sau thanh toán hợp lệ, không cộng trùng. & S1 §6 F07; S3 §4.1 Task 10 \\ \hline
Kiểm chứng và bàn giao & Kiểm thử, Docker Compose, README, cấu hình mẫu và demo; ghi lại môi trường và phạm vi kết quả. & S1 §4.1; S2 §0; S3 §4.1 Task 1, 9 \\ \hline
\end{longtable}
\endgroup

\subsubsection{Thanh toán VNPay và đối soát}

\begingroup\small
\renewcommand{\arraystretch}{1.2}
\begin{longtable}{>{\raggedright\arraybackslash}p{0.20\linewidth}>{\raggedright\arraybackslash}p{0.43\linewidth}>{\raggedright\arraybackslash}p{0.25\linewidth}}
\caption{Phạm vi thanh toán và đối soát}\label{tab:auto-context-and-scope-4}\\
\hline
\textbf{Nhóm} & \textbf{Nội dung trong phạm vi} & \textbf{Nguồn} \\ \hline
\endfirsthead
\textbf{Nhóm} & \textbf{Nội dung mở rộng kế thừa} & \textbf{Nguồn} \\ \hline
\endhead
VNPay Sandbox & Khởi tạo/redirect VNPay Sandbox mặc định; xác minh IPN/QueryDR. Không bao gồm thanh toán tiền thật hoặc kích hoạt cổng production. & S1 §3, §4.1, §9.9; S3 §1.1, §4 \\ \hline
Payment/Order nâng cao & Payment PENDING/SUCCESS/FAILED, shared finalization, retry có kiểm soát, điều chỉnh chính sách TTL/grace và xử lý hết hạn/hoàn tài nguyên khi có VNPay PENDING, reconciliation theo BR-01--BR-06. & S1 §6 F04--F06, §9.9, §16; S3 §4 \\ \hline
ZERO\_AMOUNT & Xử lý nội bộ đơn 0 VND qua finalization chung; độc lập Mock, không gọi VNPay/Mock cho đơn 0 VND. & S1 §9.9 BR-06, §16; S3 §4 \\ \hline
\end{longtable}
\endgroup

Các yêu cầu đặt hàng, thanh toán và đối soát được kiểm tra theo các lớp unit, integration và E2E được trình bày trong chương kiểm thử.
\par

\subsection{Out of Scope và Deferred}

\subsubsection{Out of Scope -- ngoài phạm vi sản phẩm baseline}

\begingroup\small
\renewcommand{\arraystretch}{1.2}
\begin{longtable}{>{\raggedright\arraybackslash}p{0.24\linewidth}>{\raggedright\arraybackslash}p{0.68\linewidth}}
\caption{Nội dung ngoài phạm vi}\label{tab:auto-context-and-scope-5}\\
\hline
\textbf{Nội dung} & \textbf{Ranh giới áp dụng} \\ \hline
\endfirsthead
\textbf{Nội dung} & \textbf{Ranh giới áp dụng} \\ \hline
\endhead
Thanh toán tiền thật và refund tự động & Không kích hoạt cổng thanh toán production hoặc giao dịch thật; không tự thêm nghiệp vụ refund. \\ \hline
Multi-vendor, microservices, mobile native & Giữ mô hình một cửa hàng và stack hiện có. \\ \hline
Seller/Admin phức tạp và UI Barista & Module Seller/Admin phức tạp và UI Barista nằm ngoài phạm vi sản phẩm baseline; vòng đời trạng thái đơn vẫn được mô tả theo yêu cầu kế thừa. \\ \hline
Vận chuyển và định vị thời gian thực & Không bổ sung Last-Mile, tracking hoặc nghiệp vụ giao nhận. \\ \hline
ERP/POS và quản lý nguyên liệu đầy đủ & Tồn kho chỉ phục vụ phạm vi chống oversell và lifecycle tài nguyên của baseline. \\ \hline
\end{longtable}
\endgroup

\subsubsection{Deferred -- chức năng chưa đưa vào phạm vi hiện hành}

\begingroup\small
\renewcommand{\arraystretch}{1.2}
\begin{longtable}{>{\raggedright\arraybackslash}p{0.24\linewidth}>{\raggedright\arraybackslash}p{0.68\linewidth}}
\caption{Chức năng Deferred}\label{tab:auto-context-and-scope-6}\\
\hline
\textbf{Nội dung} & \textbf{Giới hạn} \\ \hline
\endfirsthead
\textbf{Nội dung} & \textbf{Giới hạn} \\ \hline
\endhead
Forgot/Reset Password & Tiếp tục Deferred theo baseline; không bổ sung hoặc mô tả endpoint/UI tương ứng như đã có. [S1 §4.2] \\ \hline
Hủy đơn chủ động -- FR-ORD-005 & Tiếp tục Deferred; không mặc định có Cancel API/UI. Hết hạn đơn và hoàn tài nguyên tự động theo lifecycle vẫn nằm trong phạm vi kế thừa. [S1 §6 F04] \\ \hline
Advanced monitoring & Deferred theo nguồn; chưa có evidence triển khai advanced monitoring như một năng lực của sản phẩm. [S1 §4.2] \\ \hline
\end{longtable}
\endgroup

\subsection{Giả định và ràng buộc}

\begingroup\small
\renewcommand{\arraystretch}{1.2}
\begin{longtable}{>{\raggedright\arraybackslash}p{0.20\linewidth}>{\raggedright\arraybackslash}p{0.43\linewidth}>{\raggedright\arraybackslash}p{0.25\linewidth}}
\caption{Giả định và ràng buộc}\label{tab:auto-context-and-scope-7}\\
\hline
\textbf{Mã / Loại} & \textbf{Nội dung} & \textbf{Nguồn} \\ \hline
\endfirsthead
\textbf{Mã / Loại} & \textbf{Nội dung} & \textbf{Nguồn} \\ \hline
\endhead
A-01 -- Giả định phân tích & Quy tắc nghiệp vụ được đặc tả thống nhất; thay đổi phạm vi cần được quản lý trong backlog. & S1; S2; S4 \\ \hline
C-01 -- Ràng buộc công nghệ & Giữ Next.js, NestJS, TypeScript, PostgreSQL, Prisma và Docker; giao tiếp HTTP/JSON giữa frontend và backend. & S1 §3; S4 \\ \hline
C-02 -- Ràng buộc truy cập và nghiệp vụ & Guest có thể xem menu; tạo đơn, thanh toán và xem lịch sử cần JWT hợp lệ và kiểm tra quyền sở hữu. Backend là nguồn quyết định giá và trạng thái. & S1 §3--5 \\ \hline
C-03 -- Ràng buộc môi trường & Các môi trường phát triển, kiểm thử và triển khai sử dụng project, database và volume riêng. & S2 §0; S4 \\ \hline
A-02 -- Ràng buộc thanh toán & Thanh toán thử nghiệm dùng Sandbox hoặc Mock theo chính sách môi trường; không xử lý tiền thật. & S1; S4 \\ \hline
C-04 -- Ràng buộc đánh giá & Kết quả được diễn giải theo môi trường và phạm vi thực nghiệm. & S1; S2 \\ \hline
\end{longtable}
\endgroup

Rủi ro liên quan tính nhất quán, thanh toán bất đồng bộ và vận hành được phân tích trong chương thiết kế, kiểm thử và mục hạn chế\ref{sec:limitations}.
